% Einheit 4 – Nr. 50–57
\setcounter{aufgabe}{49}
\einheitenkopf[e4][4 Verdopplungs- und Halbwertszeit]{Einheit 4 von 5 · Verdopplungs- und Halbwertszeit}
\zweigzeile{Hier lernst du, die Zeit zu bestimmen, nach der sich ein Wert verdoppelt oder halbiert hat · neu in diesem Jahr · P10 · baut auf: Wachstumstabelle (Einheit 2), Wachstumsgleichung (Einheit 3), Werte am Graphen ablesen}

\begin{aufgabe}{Ich erkenne, ob sich die Zeit oder der Bestand verdoppeln oder halbieren soll.}
Kreuze an, was sich verdoppeln oder halbieren soll. Rechne nichts.
\begin{teile}
\teil Nach wie vielen Jahren hat sich das Guthaben verdoppelt? \quad \kreuz{die Zeit}\kreuz{der Bestand}
\teil Wie lange dauert es, bis nur noch die Hälfte des Wirkstoffs im Blut ist? \quad \kreuz{die Zeit}\kreuz{der Bestand}
\teil Wie viele Tiere sind es, wenn man doppelt so lange wartet? \quad \kreuz{die Zeit}\kreuz{der Bestand}
\teil Nach welcher Zeit sind aus 300 Bakterien 600 geworden? \quad \kreuz{die Zeit}\kreuz{der Bestand}
\teil Wie viel ist nach der halben Beobachtungszeit noch da? \quad \kreuz{die Zeit}\kreuz{der Bestand}
\end{teile}
\end{aufgabe}

\begin{aufgabe}{Ich kann den Zielwert für die Verdopplung oder Halbierung angeben.}
Gib den Wert an, der erreicht sein muss.
\begin{teile}
\teil Verdopplung, Start bei 40 \quad \leerfeld
\teil Verdopplung, Start bei 250 \quad \leerfeld
\teil Halbierung, Start bei 90 \quad \leerfeld
\teil Halbierung, Start bei 3,6\,g \quad \leerfeld[g]
\teil Ein Guthaben startet bei 500\,€ und hat nach fünf Jahren 700\,€. Welchen Betrag muss es für die Verdopplung erreichen? \leerfeld[€]
\end{teile}
\end{aufgabe}

\begin{aufgabe}{Ich kann die Verdopplungs- oder Halbwertszeit aus einer Tabelle ablesen.}
Lies die Verdopplungszeit ab.
\begin{teile}
\teil \wertetabelle[5,10,20,40]{t \text{ in h}}{N}{0,1,2,3} \quad \feld[h]{T}
\end{teile}
\anweisung{Der Zielwert liegt zwischen zwei Tabellenwerten. Gib die nächstliegende Zeit an.}
\begin{teile}
\teil Verdopplungszeit: \\ \wertetabelle[16,24,36,54,81]{t \text{ in Jahren}}{N}{0,1,2,3,4} \quad \feld[Jahre]{T}
\teil Halbwertszeit: \\ \wertetabelle[80,56,{39{,}2},{27{,}44}]{t \text{ in h}}{m \text{ in mg}}{0,1,2,3} \quad \feld[h]{T}
\teil Ein radioaktives Präparat zerfällt. Gib seine Halbwertszeit an. (P10 2016 OS) \\ \wertetabelle[12,{9{,}6},{7{,}68},{6{,}14},{4{,}92}]{t \text{ in Tagen}}{m \text{ in mg}}{0,2,4,6,8} \quad \feld[Tage]{T}
\end{teile}
\end{aufgabe}

\begin{aufgabe}{Ich kann die Verdopplungs- oder Halbwertszeit am Graphen ablesen.}
Links: Wirkstoff im Blut ($t$ in Stunden, $m$ in mg). Rechts: Zahl der Kaninchen auf einer Insel ($t$ in Monaten).

\begin{ksys}[xmin=0,xmax=8,ymin=0,ymax=400,ystep=50,ablesen,xlabel=$t$,ylabel=$m$]
\funktion{400*0.8^\x}{}
\end{ksys}\ksysabstand\begin{ksys}[xmin=0,xmax=10,ymin=0,ymax=180,ystep=20,ablesen,xlabel=$t$,ylabel=$N$]
\funktion{60*1.12^\x}{}
\end{ksys}

\begin{teile}
\teil Links: Lies ab, nach welcher Zeit nur noch die Hälfte des Wirkstoffs im Blut ist. \feld[h]{T}
\teil Beschreibe in zwei Sätzen, wie du bei a) vorgegangen bist.
\teil Rechts: Ermittle mithilfe des Graphen die Verdopplungszeit der Kaninchenzahl und beschreibe dein Vorgehen. (P10 2020 OS)
\end{teile}
\end{aufgabe}

\begin{aufgabe}{Ich kann die Verdopplungs- oder Halbwertszeit durch Probieren mit dem Faktor bestimmen.}
Bestimme durch Probieren, nach wie vielen ganzen Schritten sich der Wert ungefähr verdoppelt oder halbiert hat.
\begin{teile}
\teil Faktor 1,5, Verdopplung \quad \leerfeld
\teil Faktor 1,25, Verdopplung \quad \leerfeld
\teil Faktor 0,5, Halbierung \quad \leerfeld
\teil Faktor 0,8, Halbierung \quad \leerfeld
\teil Zunahme um 10\,\% je Schritt, Verdopplung \quad \leerfeld
\teil Abnahme um 10\,\% je Schritt, Halbierung \quad \leerfeld
\anweisung{Prüfe mit einer Rechnung.}
\teil Ein Wert wächst jährlich um 10\,\%. Berechne für die Startwerte 80 und 500 den Wert nach sieben Jahren. Was fällt dir auf?
\end{teile}
\anweisung{Löse mit dem Logarithmus (Taste log).}
\begin{teile}
\teil Bestimme die Verdopplungszeit bei 10\,\% Zunahme pro Jahr auf zwei Nachkommastellen genau. \feld[Jahre]{T}
\end{teile}
\end{aufgabe}

\begin{aufgabe}{Ich finde den Fehler bei der Verdopplungs- oder Halbwertszeit.}
Finde den Fehler, benenne ihn und korrigiere.
\begin{teile}
\teil Ein Bestand wächst. Paul sagt: „Die Verdopplungszeit ist 2 Jahre, denn 2 Jahre sind das Doppelte von 1 Jahr.“ \\ \wertetabelle[300,360,432,{518{,}4},{622{,}08}]{t \text{ in Jahren}}{N}{0,1,2,3,4}
\teil Lea soll am Graphen die Halbwertszeit eines Stoffes ablesen, der mit 8\,mg beginnt. Sie sucht die 4 auf der waagerechten Achse, geht senkrecht bis zum Graphen und liest an der senkrechten Achse den Wert ab.
\teil Ein Stoff zerfällt. Maja sagt: „Die Halbwertszeit ist 3 Stunden, weil 29,64 der letzte Wert über 25 ist.“ \\ \wertetabelle[50,42,{35{,}28},{29{,}64},{24{,}89}]{t \text{ in h}}{m \text{ in mg}}{0,1,2,3,4}
\end{teile}
\end{aufgabe}

\begin{aufgabe}{Ich kann begründen, wann es eine feste Verdopplungszeit gibt.}
Begründe.
\begin{teile}
\teil Ein Wert wächst jedes Jahr um 10\,\%. Warum dauert es von 100 bis 200 genauso lange wie von 1000 bis 2000?
\teil Ein Wert wächst linear: 100; \ 120; \ 140; \ 160; \ … \ Gibt es eine feste Verdopplungszeit?
\end{teile}
\end{aufgabe}

\begin{aufgabe}{Ich kann mit der Halbwertszeit eine Frage im Sachzusammenhang beantworten.}
Ein Schmerzmittel hat im Blut eine Halbwertszeit von 6 Stunden. Um 8 Uhr sind 240\,mg im Blut.
\begin{teile}
\teil Wie viel ist um 14 Uhr, um 20 Uhr und um 2 Uhr nachts noch im Blut?
\teil Eine neue Tablette darf erst genommen werden, wenn weniger als 70\,mg im Blut sind. Zu welcher der Uhrzeiten aus a) ist das zuerst der Fall?
\teil Eine Freundin sagt: „Nach 24 Stunden ist nichts mehr im Blut.“ Stimmt das? Begründe mit einer Rechnung.
\end{teile}
\end{aufgabe}
